\scanpage{57}
\mspage{9}
Notons que l'isom
\[
\psi_r:\Pi_{03}\xrightarrow{\sim}\Pi_\omega
\]
\inserted{$r\in\operatorname{Rep}(J,\uline{\unsure{\bar\omega|\varpi}})\ \deleted{\unsure{pour}}\simeq\operatorname{Isom}((\Sigma_0^*,P_+),(\Sigma^*,P_+))$}
est compatible avec le transport des structures, \deleted{ce qui fait} \inserted{donc pour} $\forall$ élément
\[
g\in\widetilde{\mathcal G}_J\subset\Gamma=\mathcal G_J
\]
\inserted{
\[
\begin{aligned}
\widetilde{\mathcal G}_J=\Gamma_{P_\varepsilon}&=\left\{(g,\tau^{\unsure{\operatorname{sg}g|\alpha}})\in\mathcal G_J\times\deleted{\{1,\tau\}}\inserted{\mathbb Z/2\mathbb Z}\ \middle|\ \begin{array}{l}\alpha=0\text{ si }\operatorname{sg}g=1\\\alpha=1\text{ si }\operatorname{sg}g=-1\end{array}\right\}\\
&=\{g\in\Gamma\mid \unsure{g_\varepsilon^{(\varepsilon)}|g(\varepsilon)}=\operatorname{id}_{\unsure{\varepsilon}}\}
\end{aligned}
\]
}
on aura commutativité de
\diagram{
\begin{tikzcd}[ampersand replacement=\&]
\unsure{\checkmark|\vee}\&\Pi_{03}\arrow[dl,"\psi_r=\psi_{\omega,i}^\varepsilon"']\arrow[dr,"\psi_{g(r)}=\psi_{g(\omega),g(i)}^{g(\varepsilon)}"]\&\\
(37)\quad\Pi_\varepsilon\arrow[rr,"\Pi_1(g)","\text{(par transport de structure)}"']\&\&\Pi_\varepsilon
\end{tikzcd}
}
ce qui suffit en principe à passer d'un $\psi_r$ à $\psi_{g(r)}=\deleted{\unsure{\varphi}}\Pi_1(g)\circ\psi_r$, vu que \deleted{\unsure{g}} $\widetilde{\mathcal G}_J$ est simplement transitif sur l'ens. des repères $r$ correspondant à un $\varepsilon$ fixé. Mais, on aimerait exprimer le passage par une formule de transformation \uline{dans} $\Pi_{0,3}$.

Notons que la donnée du repère \inserted{$r$} définit un isom canonique (encore noté $u_r$)
\margin{56}
